\section{Animal Classifier}
\label{sec:animal_classifier}

i\subsection*{Learnings \& Troubleshooting}

\textbf{Problem: Falsche Bildauflösung (Resolution Shift)} \\
Wenn ein Convolutional Neural Network (wie \texttt{efficientnet\_v2\_s}) auf kleinen $224 \times 224$~Pixel großen Bildausschnitten trainiert wird, lernt es Filter für Objekte in einer bestimmten relativen Größe. Werden bei der Evaluierung oder bei der Visualisierung mittels Grad-CAM plötzlich $384 \times 384$~Bilder übergeben, kommt es zu einem sogenannten \emph{Resolution Shift}. Die gelernten Filter passen nicht mehr optimal zu den Strukturen, was die Accuracy künstlich verschlechtern kann.

\paragraph{Empfehlung:} 
Passe die Transformationen in der Datei \texttt{animalClassifier.py} so an, dass sie exakt der Trainingsauflösung von $224 \times 224$~Pixeln entsprechen. Alternativ kann in \texttt{train.py} direkt auf $384 \times 384$~Pixel trainiert werden, was allerdings mehr GPU-Speicher benötigt.